\chapter{Гальма}
% full version: chapters/03_gaba.tex

Додамо в машину другий за чисельністю тип нейронів --- \nt{GABA}. Їхній сигнал не підштовхує сусіда до клацання, а відштовхує. Це гальма. Їх небагато, у корі від п'ятої частини до чверті всіх нейронів, але без них машина не їздить.

Є одне обмеження: «плюсовий» нейрон не може сам загальмувати сусіда. Якщо потрібен від'ємний зв'язок, його доводиться будувати через посередника: «плюс» будить гальмівний нейрон, а той уже гальмує ціль. Зміна знака коштує одного нейрона й однієї невеликої затримки.

\section*{«Але не якщо»}

Найкорисніша операція з гальмами --- заборона: «спрацюй від A, але не якщо прийшло B». Нейрон отримує поштовх від A і гальмо від посередника, якого будить B. Із заборон і «або» можна зібрати будь-яку логіку.

\pencilfig{3}{\pencilart{3}}{«А, але не якщо Б»: червоний нейрон перекриває дорогу синьому.}

Заборона в часі дає детектор руху, який є у вас в оці. Два сусідні датчики: один збуджує вихідний нейрон, другий гальмує його із затримкою. Якщо пляма світла біжить в один бік із відповідною швидкістю, гальмування приходить саме тоді, коли й збудження, --- і гасить його. Якщо в інший бік --- гальмування запізнюється, і нейрон встигає клацнути. Виходить нейрон, який бачить рух лише в один бік.

\pencilfig{103}{%
% sensors on top; the left one excites the output, the right one inhibits it through a GABA go-between and a delay
\draw[dotpen=pcAmber, wash=pcAmber] (0,1.6) circle (0.16); \draw[dotpen=pcAmber, wash=pcAmber] (1.6,1.6) circle (0.16);
\node[lbl, anchor=south] at (0.8,1.8) {два датчики};
\draw[pen=pcBlue, wash=pcBlue] (0.4,0) circle (0.24);
\draw[arr=pcBlue] (0,1.44) -- (0.3,0.25);
\draw[pen=pcBlue] (1.6,1.44) -- (1.6,1.18);
\draw[penlite, fill=white] (0.95,0.9) rectangle (2.25,1.18); \node[font=\tiny\itshape, text=pcGray] at (1.6,1.04) {затримка};
\draw[pen=pcBlue] (1.6,0.9) -- (1.6,0.72);
\draw[pen=pcRed, wash=pcRed] (1.6,0.55) circle (0.17);
\draw[pcRed, line width=0.7pt, -{Bar[width=5pt]}] (1.45,0.45) -- (0.66,0.08);
\draw[arr] (0.4,-0.24) -- (0.4,-0.6); \node[lbl, anchor=north] at (0.4,-0.6) {вихід};
% outcomes
\begin{scope}[xshift=3.1cm]
  \draw[dotpen=pcAmber, wash=pcAmber] (0,1.25) circle (0.1); \draw[arr=pcAmber] (0.18,1.25) -- (1.1,1.25);
  \node[lbl, anchor=west] at (1.25,1.25) {зліва направо: клацання};
  \draw[dotpen=pcAmber, wash=pcAmber] (1.1,0.45) circle (0.1); \draw[arr=pcAmber] (0.92,0.45) -- (0,0.45);
  \node[lbl, anchor=west] at (1.25,0.45) {справа наліво: тиша};
  \node[lbl, anchor=west, align=left] at (0,-0.35) {гальмування встигає\\якраз до збудження};
\end{scope}
}{Детектор руху: гальмування приходить із затримкою і гасить відповідь, лише якщо пляма біжить справа наліво.}


\section*{Віднімання чи ділення}

Гальмування вміє не лише віднімати. Відкритий гальмівний канал не тягне напругу вниз, він тягне її \emph{до спокою} --- ніби у відрі проробили ще одну дірку. Через це будь-який приплив води піднімає її менше. Гальмування не забирає від сигналу сталу порцію, а \emph{ділить} його: це регулятор гучності. Якщо сила гальмування зростає разом із загальною активністю навколо, кожен нейрон реагує не на те, наскільки сильний його сигнал, а на те, наскільки він сильний \emph{порівняно із сусідами}. Тому та сама мережа працює і в сутінках, і на яскравому сонці. Щоправда, на частоту імпульсів такий «дільник» діє чисто лише разом зі сталим фоновим шумом, який у живому мозку є завжди.

\pencilfig{104}{%
\foreach \dx/\t in {0/віднімання,4.6/ділення}{
  \begin{scope}[xshift=\dx cm]
  \draw[pen] (0,0) -- (3.6,0); \draw[pen] (0,0) -- (0,1.9);
  \node[lbl, anchor=north] at (1.8,-0.05) {сигнал}; \node[lbl, anchor=south, rotate=90] at (-0.05,0.95) {відповідь};
  \draw[pen=pcBlue] (0.4,0) -- (3.3,1.75);
  \node[font=\small\itshape, text=pcGray, anchor=south] at (1.8,1.95) {\t};
  \end{scope}}
\draw[pen=pcRed, line width=0.9pt] (1.3,0) -- (3.4,1.27);
\node[lbl, text=pcRed, anchor=west] at (2.25,0.35) {зсув};
\begin{scope}[xshift=4.6cm]
\draw[pen=pcRed, line width=0.9pt] (0.4,0) -- (3.3,0.8);
\node[lbl, text=pcRed, anchor=west] at (2.0,0.25) {нахил менший};
\end{scope}
\node[lbl, text=pcBlue, anchor=east] at (2.25,1.45) {без гальм};
}{Синім --- відповідь нейрона без гальмування, червоним --- із ним. Віднімання зсуває поріг; ділення зменшує гучність, як ручка регулятора.}

Є й другий наслідок: зайва дірка змушує відро швидше забувати. Вікно, у яке два сигнали мають потрапити, щоб додатися, звужується. Гальмування робить детектори збігів суворішими.

\section*{Вибір}

Тепер головне, чого бракувало: вибір одного з багатьох. Дві групи нейронів гальмують одна одну. Поки гальмування слабке, працюють обидві, просто тихіше. Але якщо воно досить сильне, будь-яка маленька різниця наростає, як на гойдалці, доки одна група не замовкне зовсім. Переможець забирає все. І перемога тримається: щоб мережа передумала, новий аргумент має бути помітно сильнішим за старий. Машина не смикається від кожного шереху.

Гальма стирають і пам'ять: гальмівний нейрон за сигналом «скидання» гасить вогнище з попереднього розділу. Дві такі комірки, що гальмують одна одну, --- аналог засувки, основи будь-якої комп'ютерної пам'яті.

\section*{Життя на межі}

Мережа з самих «плюсів» може піти в лавину: кожен імпульс породжує більше ніж один наступний. Гальма тримають її в робочій точці, як регулятор тримає температуру в духовці. Умов дві: гальмування має бути досить сильним і досить швидким. Якщо воно сильне, але повільне, регулятор починає перестаратися --- як людина, що крутить кран у душі, не дочекавшись, поки дійде гаряча вода. Температура гойдається туди-сюди. У мозку така петля з «плюсів» і «мінусів» породжує швидкий ритм --- кілька десятків коливань на секунду.

Кора, схоже, живе саме на межі: її власний зворотний зв'язок досить сильний, щоб без гальм піти в лавину. Утримують її лише швидкі гальма, і тому вона так чутливо відповідає на слабкі сигнали.

Висновок розділу несподіваний: гальма майже не додають машині нових обчислень. Вони додають \emph{керування}: вибір, скидання, регулювання гучності, стійкість. Збудження несе дані, гальмування вирішує, яким даним пройти, коли і наскільки гучно.

\fullversion{3}
